\documentclass[10pt,a4paper]{amsart}
\usepackage[margin=19mm]{geometry}
\usepackage{amsmath,amssymb,mathtools}
\usepackage[scr=rsfs,cal=euler]{mathalfa}
\usepackage{xfrac}
\usepackage[unicode,hidelinks,bookmarks=false]{hyperref}
\newcommand{\ie}{i.e.\ }
\newcommand{\cf}{cf.\ }
\newcommand{\resp}{respectively\ }
\newcommand{\Subsection}{\subsection}
\providecommand{\nobreakdash}{\nobreak\hbox{-}}
\setlength{\emergencystretch}{2em}
\multlinegap0pt
\usepackage{bm}
\usepackage{bbm}
\usepackage{dsfont}
\usepackage{xspace}
\usepackage{cancel}
\usepackage{mathtools}
\usepackage{amsthm}
\makeatletter
\@ifundefined{theorem}{\newtheorem{theorem}{Theorem}}{}
\@ifundefined{lemma}{\newtheorem{lemma}[theorem]{Lemma}}{}
\makeatother
\newcommand{\eqdef }{\overset{\mbox{\tiny{def}}}{=}}
\title[Local Existence and Finite-Time Singularity Formation in the]{Local Existence and Finite-Time Singularity Formation in the Vlasov-Poisson-Isotropic Landau System}
\author{Jin Woo Jang, Junsung Kim}
\date{Source: arXiv:2603.14871v1 (CC BY 4.0)}
\begin{document}
\maketitle

\noindent\emph{Source and licence.} Local Existence and Finite-Time Singularity Formation in the Vlasov-Poisson-Isotropic Landau System. Version arXiv:2603.14871v1 (CC BY 4.0), distributed under the Creative Commons Attribution 4.0 International licence, as recorded in the arXiv licence metadata for this exact version and shown on the abstract page https://arxiv.org/abs/2603.14871v1. Full rights notice: licenses/arxiv-2603.14871-CC-BY-4.0.txt.\par\smallskip

\noindent\textit{Editorial scope.} Counted unit: the Lemma whose source label is \texttt{local comparison} in the article (source0.tex lines 1086--1093, source label \texttt{local comparison}), delivered together with the article's own complete original proof of it (source0.tex lines 1095--1256, one \texttt{proof} environment of 162 lines). No mathematics is written, repaired or re-derived: statement and proof are the article's own text line for line. Required context retained uncounted: def.Leg (lines 779--781); def cutoff (lines 785--787); reg interval (lines 788--790); semi (lines 793--795); approx initial (lines 800--802); int decay (lines 831--833); semi problem (lines 1003--1005); f equals fin (lines 1042--1044); embed (lines 1060--1064). Every definition, display and label of those lines is retained unchanged. Omitted from this package and not redistributed: 1--778, 782--784, 791--792, 796--799, 803--830, 834--1002, 1006--1041, 1045--1059, 1065--1085, 1094--1094, 1257--1994. Nothing inside a retained excerpt is omitted or altered: every retained body is delivered exactly as the article's lines are written, so no omission receipt of any kind accompanies this package. Citations: the retained excerpts carry no citation command at all, so no bibliography entry of the article is transcribed and no reference is redistributed. Reference adaptation: none was needed and none was made; every reference inside the delivered unit resolves inside it, so no unresolved reference or citation is produced. Printed slips kept literal: no printed word, symbol, hypothesis or number of the counted statement or of its proof was altered. Layout and class: the article's own class is \texttt{amsart} with packages including showlabels, geometry, amsaddr, amssymb, verbatim, amsthm, amsmath, amsfonts, latexsym, mathtools, changepage, enumitem, babel, esint; the wrapper is the lane's pinned \texttt{amsart} editorial wrapper at 10pt on A4, which restates the theorem-like environments the retained text needs and adds the article's own macro definitions that the retained text needs (\texttt{\textbackslash eqdef}), with extra wrapper packages (bm, bbm, dsfont, xspace, cancel, mathtools). The delivered numbering is the unit's own. Assets: no retained span contains a figure environment, a graphics-inclusion command, a PDF-page inclusion, a table or a TikZ picture; no image file is copied into this package, the package declares no asset dependency and no image file is redistributed with it. Licence: version arXiv:2603.14871v1 of this article is distributed under the Creative Commons Attribution 4.0 International licence, as recorded in the arXiv licence metadata for this exact version and shown on the abstract page https://arxiv.org/abs/2603.14871v1. A full rights notice is delivered at licenses/arxiv-2603.14871-CC-BY-4.0.txt. Characters: every retained excerpt is pure ASCII, so the delivered engine, XeLaTeX with its default Latin Modern text font, reports no missing character.
\begin{align}\label{def.Leg}
\mathcal{L}_{\varepsilon,g}(f) \eqdef  \partial_t f - \varepsilon \Delta_{x,v}f + \chi(\varepsilon v)v\cdot\nabla_x f + J^{\varepsilon}(t)\chi(\varepsilon x) E_g\cdot\nabla_v f - (-\Delta_v)_\varepsilon^{-1}g \Delta_vf - gf,
\end{align}where we define the localized diffusion coefficient

\begin{align}\label{def cutoff}
0\le \chi\le 1,\quad \chi\equiv1\,\,\textup{ in }\,\,B_{\frac{1}{2}}(0), \quad \textup{ and }\quad \textup{supp }\chi=B_1(0),
\end{align} and $J^\varepsilon$ is a time-regularized characteristic function defined via the following mollification:

\begin{align}\label{reg interval}
J^\varepsilon(t) &\eqdef \left(1_{\left[-\frac{\varepsilon}{2},T-\frac{\varepsilon}{2}\right]}\star \rho_\varepsilon\right)(t),
\end{align}with a mollifier $\rho_\varepsilon$ supported on $\left[-\frac{\varepsilon}{2}, \frac{\varepsilon}{2}\right]$.

\begin{align}\label{semi}
\mathcal{L}_{\varepsilon,g}(f) = \mathcal{L}_{\varepsilon,g}(f^\textup{in}_\varepsilon)I^\varepsilon,
\end{align}where

\begin{align}\label{approx initial}
    \{f^\textup{in}_\varepsilon\ge0\}_{\varepsilon>0}\subset C^\infty_c(\mathbb{R}^6)\textup{ s.t. } f^\textup{in}_\varepsilon \to f^\textup{in} \textup{ in }\mathcal{C}^2_{m,0}(\mathbb{R}^6),
\end{align}where the approximates are just obtained from multiplying $f^\textup{in}$ by a $6$-dimensional smooth cut-off. Here, another time-regularized characteristic function near $t=0$, $I^\varepsilon$ is defined by

\begin{equation}\label{int decay}
    \int_{\mathbb{R}^3}\frac{\langle z\rangle^{-m}}{|x-z|^k}dz \lesssim_{m,k} \frac{1}{\langle x\rangle^k}, \text{ for }m>3\text{ and }0<k<3.
\end{equation}

\begin{align}\label{semi problem}
\mathcal{L}_{\varepsilon,g}(f) = \mathcal{L}_{\varepsilon,g}(f^\textup{in}_\varepsilon)I^\varepsilon \quad \textup{ in }\,\, \mathcal{Q}_T,\textup{ and }f(t,\cdot,\cdot) = f^\textup{in}_\varepsilon(\cdot,\cdot) \quad \textup{ in }\,\, t\in[0,\varepsilon),
\end{align} where $\mathcal{L}_{\varepsilon,g}$ is defined as in \eqref{def.Leg}.

\begin{equation}\label{f equals fin}
f(t,x,v) = f^\textup{in}_\varepsilon(x,v) \quad \textup{ for }\,\,t\in[0,\varepsilon).
\end{equation}This completes the construction.

\begin{lemma}\label{embed}Let $\varepsilon>0, T>0$, and $m>0$. Then for all $(t,x,v)\in\mathcal{Q}_T$, we have
\begin{align*}
 \langle x\rangle^{-m} \langle v\rangle^{-m}\le (2\langle T\rangle)^m \langle x-\chi(\varepsilon v)vt\rangle^{-m},
\end{align*}where $\chi(\cdot)$ is the cut-off function defined as in \eqref{def cutoff}.
\end{lemma}

\begin{lemma}[Function-wise upper and lower bound]\label{local comparison} Let $g \in Y_{T,\alpha,p_0,q_0,w}$ satisfy $g\ge0$ and
\begin{align*}
 g(t,x,v) \le C_g\langle x-\chi(\varepsilon' v)vt \rangle^{-m}\langle v\rangle^{-m},
\end{align*}for some $C_g\ge 0$ and $0<\varepsilon'\le 1$. Let $f\in Y_{T,\alpha,p_0,q_0,w}$ be a solution to the linear problem \eqref{semi problem} for some $0<\varepsilon\le \varepsilon'$ with the approximated initial data $f^\textup{in}_\varepsilon\in C^\infty_c(\mathbb{R}^6)$ (see the definition \eqref{approx initial}) satisfying $f^\textup{in}_\varepsilon\ge0$. Then, there exists a constant $M=M(m)>0$ such that
\begin{align*}
 f^\textup{in}_\varepsilon I^\varepsilon\le f(t,x,v) \le e^{M(T+1)^3(C_g+1)}\lVert f^\textup{in}_\varepsilon\rVert_{\mathcal{C}^2_{m,0}}(2\langle T\rangle)^m\langle x -\chi(\varepsilon v)vt\rangle^{-m} \langle v\rangle^{-m},
\end{align*} for $(t,x,v)\in\overline{\mathcal{Q}_T}$ where $I^\varepsilon$ is defined as in \eqref{reg interval}.
\end{lemma}

\begin{proof}
\textbf{Upper bound on $t\in[0,\varepsilon]$.} In this interval, we have $f(t)=f^\textup{in}_\varepsilon$ by \eqref{f equals fin}. Then by Lemma \ref{embed}, we have 
\begin{align*}
|f(t,x,v)| \le \langle x\rangle^{-m} \langle v\rangle ^{-2m} |\langle x\rangle^m \langle v\rangle ^{2m} f^\textup{in}_\varepsilon|\le \lVert f^\textup{in}_\varepsilon\rVert_{\mathcal{C}^2_{m,0}}(2\langle T\rangle)^m\langle x-\chi(\varepsilon v)vt\rangle^{-m}\langle v\rangle^{-m}.
\end{align*}This proves the desired upper bound with $M=0$.


\noindent\textbf{Upper bound on }$t\in[\varepsilon,T]$\textbf{.} We define the auxiliary function
\begin{align*}
\overline{f}_\varepsilon(t,x,v) = e^{C'(t-\varepsilon)}(2\langle T\rangle)^m\lVert f^\textup{in}_\varepsilon\rVert_{\mathcal{C}^2_{m,0}(\mathbb{R}^6)} \langle x - \chi(\varepsilon v)vt\rangle^{-m}\langle v\rangle^{-m},
\end{align*}with a sufficiently large $C'>0$ which will be determined as in \eqref{final C'}. Note that by Lemma \ref{embed} again, we have $\overline{f}_\varepsilon(\varepsilon,x,v)\ge f^\textup{in}_\varepsilon(x,v)= f(\varepsilon, x,v)$. Therefore, by the comparison principle, to establish $f\le \overline{f}_\varepsilon$, it suffices to verify the following:
\begin{align}\label{4.2 desired ineq}
\mathcal{L}_{\varepsilon,g}(\overline{f}_\varepsilon) \ge \mathcal{L}_{\varepsilon,g}(f)=\mathcal{L}_{\varepsilon,g}(f^\textup{in}_\varepsilon)I^\varepsilon  \quad \textup{ in }\,\, (\varepsilon, T)\times\mathbb{R}^6,
\end{align}where $\mathcal{L}_{\varepsilon,g}$ denotes the linear differential operator considered in \eqref{semi}.

To this end, we proceed to compute and estimate each term in $\mathcal{L}_{\varepsilon,g}(\overline{f}_\varepsilon)$. 

\noindent \textit{(I) Transport term:} Note that
\begin{align}\label{Trans est}
(\partial_t + \chi(\varepsilon v)v\cdot\nabla_x)\overline{f}_\varepsilon = C'\overline{f}_\varepsilon,
\end{align}since $(\partial_t + \chi(\varepsilon v) v\cdot\nabla_x)\langle x-\chi(\varepsilon v)vt\rangle^{-m} \equiv 0$.

\noindent \textit{(II) Acceleration term:} To estimate the acceleration term $J^\varepsilon(t)\chi(\varepsilon x)E_g\cdot\nabla_v\overline{f}_\varepsilon$, we begin with computing the velocity gradient 
\begin{align*}
\nabla_v\langle x-\chi(\varepsilon v)vt\rangle^{-m}.
\end{align*}Using the chain rule, this term can be written as
\begin{multline*}
\nabla_v\langle \chi(\varepsilon v)vt - x\rangle^{-m} = -\frac{m}{2} \langle \chi(\varepsilon v)vt - x\rangle^{-m-2} \nabla_v|\chi(\varepsilon v)vt - x|^2\\
= -m t\langle \chi(\varepsilon v)vt - x\rangle^{-m-2}\left[\chi(\varepsilon v)(\chi(\varepsilon v)vt - x) + \varepsilon(\nabla_v\chi)(\varepsilon v)\left\{v\cdot(\chi(\varepsilon v)vt - x)\right\}\right].
\end{multline*}Now apply the bounds, $0\le \chi\le1$, $|\nabla_v\chi|<\infty$ and $|v|\le\langle v\rangle$, we obtain the estimate:
\begin{align}\label{v grad}
|\nabla_v\langle x-\chi(\varepsilon v)vt\rangle^{-m}| \le C(m,\chi)T\langle x-\chi(\varepsilon v)vt\rangle^{-m}.
\end{align}
To estimate the electric field $E_g$, we use the bound:
\begin{align*}
|E_g(t,x)| \le C_g\sup_{(t,x)\in[0,T]\times\mathbb{R}^3} \int_{\mathbb{R}^3} \frac{dv}{\langle v\rangle^m} \int_{\mathbb{R}^3} \frac{dy}{|x-\chi(\varepsilon v)vt-y|^2\langle y \rangle^m}.
\end{align*}By using \eqref{int decay}, we conclude 
\begin{align}\label{E bound}
\lVert E_g\rVert_{L^\infty([0,T]\times\mathbb{R}^3)} \lesssim_m C_g.
\end{align}
Therefore, we have
\begin{align}
|J^\varepsilon(t)\chi(\varepsilon x)E_g \cdot\nabla_v\overline{f}_\varepsilon| &\le \lVert E_g\rVert_{L^\infty_{t,x}} |\nabla_v\overline{f}_\varepsilon| \le C(m,\chi)\lVert E_g\rVert_{L^\infty_{t,x}} (T+1)\overline{f}_\varepsilon\nonumber\\
&\le C_1(m,\chi)(T+1)C_g\overline{f}_\varepsilon.\label{Acc est}
\end{align}Here, the constant $C_1$ may be redefined, if necessary, to absorb numerical factors.

\noindent \textit{(III) Nonlocal collision term:} To estimate the term $(-\Delta_v)^{-1}g\Delta_v\overline{f}_\varepsilon$, we first compute the velocity Laplacian of the weight function:
\begin{align*}
\Delta_v\langle x-\chi(\varepsilon v)vt\rangle^{-m}.
\end{align*}By direct computation, this can be written as:
\begin{multline*}
\Delta_v\langle x-\chi(\varepsilon v)vt\rangle^{-m} 
= -mt\left[\nabla_v\langle x-\chi(\varepsilon v)vt\rangle^{-m-2}\right]\cdot [\chi(\varepsilon v)(\chi(\varepsilon v)vt - x) \\
+\varepsilon(\nabla_v\chi)(\varepsilon v) \left\{v\cdot(\chi(\varepsilon v)vt -x)\right\} ] \\
- mt \langle \chi(\varepsilon v)vt - x\rangle^{-m-2}\nabla_v\cdot\left[\chi(\varepsilon v)(\chi(\varepsilon v)vt - x) + \varepsilon(\nabla_v\chi)(\varepsilon v) \left\{v\cdot(\chi(\varepsilon v)vt -x)\right\}\right]\\
=: \Delta_1 + \Delta_2.
\end{multline*}We first estimate $\Delta_1$. Noting that
\begin{multline*}
\nabla_v\langle x-\chi(\varepsilon v)vt\rangle^{-m-2} = -(m+2)t \langle \chi(\varepsilon v) vt - x\rangle^{-m-4}[\chi(\varepsilon v)(\chi(\varepsilon v)vt - x) \\
+ \varepsilon(\nabla_v\chi)(\varepsilon v)\left\{v\cdot(\chi(\varepsilon v)vt - x)\right\}],
\end{multline*}we deduce the bound
\begin{align}\label{Delta 1}
|\nabla_v\langle x-\chi(\varepsilon v)vt\rangle^{-m-2}| \le (1+\|\nabla_v\chi\|_{L^\infty})(m+2)T\langle \chi(\varepsilon v)vt - x\rangle^{-m-3}.
\end{align}Next, we simplify and estimate $\Delta_2$. For the first divergence component:
\begin{multline*}
\nabla_v\cdot\left[\chi(\varepsilon v)(\chi(\varepsilon v)vt - x)\right]
=\nabla_v\left[\chi(\varepsilon v)\right]\cdot(\chi(\varepsilon v)vt - x) + \chi(\varepsilon v)\nabla_v\cdot(\chi(\varepsilon v) vt - x)\\
= \varepsilon(\nabla_v\chi)(\varepsilon v)\cdot (\chi(\varepsilon v)vt - x) + \chi(\varepsilon v)\left[3t\chi(\varepsilon v) + v\cdot\varepsilon (\nabla_v\chi)(\varepsilon v)t\right],
\end{multline*}yielding the bound
\begin{align}\label{Delta 2-1}
|\nabla_v\cdot[\chi(\varepsilon v)(\chi(\varepsilon v)vt - x)]| \le \|\nabla_v\chi\|_{L^\infty}|\chi(\varepsilon v)vt - x| + (3+\|\nabla_v\chi\|_{L^\infty})T.
\end{align}For the second divergence component, we compute:
\begin{multline*}
\nabla_v\cdot \left[\varepsilon (\nabla_v\chi)(\varepsilon v)\left\{v\cdot(\chi(\varepsilon v)vt - x)\right\}\right]\\
=\varepsilon^2(\Delta_v\chi)(\varepsilon v)\left\{v\cdot(\chi(\varepsilon v)vt - x)\right\} + \varepsilon (\nabla_v\chi)(\varepsilon v)\cdot\nabla_v\left\{v\cdot(\chi(\varepsilon v)vt - x) \right\}\\
=\varepsilon^2(\Delta_v\chi)(\varepsilon v)\left\{v\cdot(\chi(\varepsilon v)vt - x)\right\} 
\\
+ \varepsilon(\nabla_v\chi)(\varepsilon v)\cdot\left[3(\chi(\varepsilon v)vt - x) + tv(\varepsilon(\nabla_v\chi)(\varepsilon v)\cdot v) + 3tv\chi(\varepsilon v)\right].
\end{multline*}Then using $0<\varepsilon\le 1$, we obtain
\begin{multline}
    |\nabla_v\cdot \left[\varepsilon (\nabla_v\chi)(\varepsilon v)\left\{v\cdot(\chi(\varepsilon v)vt - x)\right\}\right]| \\
    \le\|\Delta_v\chi\|_{L^\infty}|\chi(\varepsilon v)vt - x| + 3\|\nabla_v\chi\|_{L^\infty}|\chi(\varepsilon v) vt - x| + \|\nabla_v\chi\|_{L^\infty}^2T + 3\|\nabla_v\chi\|_{L^\infty}T\\
= (\|\Delta_v\chi\|_{L^\infty}+ 3\|\nabla_v\chi\|_{L^\infty})|\chi(\varepsilon v) vt-x| + (\|\nabla_v\chi\|_{L^\infty}^2 + 3\|\nabla_v\chi\|_{L^\infty})T.\label{Delta 2-2}
\end{multline}Combining these estimates \eqref{Delta 1}, \eqref{Delta 2-1} and \eqref{Delta 2-2}, we estimate the full Laplacian:
\begin{multline}
|\Delta_v\langle \chi(\varepsilon v) vt - x\rangle^{-m}|\le |\Delta_1| + |\Delta_2| \\
\le mT  (1+\|\nabla_v\chi\|_{L^\infty})^2(m+2)T\langle \chi(\varepsilon v)vt - x\rangle^{-m-3} |\chi(\varepsilon v)vt - x|\\
+ mT \langle \chi(\varepsilon v)vt - x\rangle^{-m-2} (\|\nabla_v\chi\|_{L^\infty}|\chi(\varepsilon v)vt - x| + (3+\|\nabla_v\chi\|_{L^\infty})T \\+ (\|\Delta_v\chi\|_{L^\infty}+ 3\|\nabla_v\chi\|_{L^\infty})|\chi(\varepsilon v)vt - x| + (\|\nabla_v\chi\|_{L^\infty}^2 + 3\|\nabla_v\chi\|_{L^\infty})T)\\
\le C(m,\chi)(T+1)^2 \langle \chi(\varepsilon v)vt - x\rangle^{-m}.\label{Delta}
\end{multline}
To complete the estimate, we compute the Laplacian of the full weight function:
\begin{multline*}
\Delta_v\left(\langle x-\chi(\varepsilon v)vt\rangle^{-m} \langle v\rangle^{-m}\right) = \Delta_v\left(\langle x-\chi(\varepsilon v)vt\rangle^{-m}\right)\langle v\rangle^{-m} \\ +2\nabla_v\langle x-\chi(\varepsilon v)vt\rangle^{-m}\cdot\nabla_v\langle v\rangle^{-m} +
\langle x-\chi(\varepsilon v)vt\rangle^{-m}\Delta_v\left(\langle v\rangle^{-m}\right).
\end{multline*}Applying the previous bounds, \eqref{v grad} and \eqref{Delta}, we obtain
\begin{align}\label{total Delta}
 \Delta_v\left(\langle x-\chi(\varepsilon v)vt\rangle^{-m} \langle v\rangle^{-m}\right)\le C(m,\chi)(T+1)^2\langle x-\chi(\varepsilon v)vt\rangle^{-m}\langle v\rangle^{-m}.
\end{align}

Altogether, we obtain the following upper bound for the localized diffusion term:
\begin{align*}
|(-\Delta_v)_\varepsilon^{-1}g\Delta_v\overline{f}_\varepsilon| \le C_2(m,\chi)(T+1)^2\lVert (-\Delta_v)_\varepsilon^{-1}g\rVert_{L^\infty_{t,x,v}} \overline{f}_\varepsilon.
\end{align*}Moreover, since
\begin{align*}
(-\Delta_v)_\varepsilon^{-1}g \le C_g \int_{\mathbb{R}^3} \frac{du}{|u-v|\langle u\rangle^m} \le C_2(m)C_g,
\end{align*}we conclude that
\begin{align}\label{Coll1 est}
|(-\Delta_v)_\varepsilon^{-1}g\Delta_v\overline{f}_\varepsilon| \le C_2(m,\chi)(T+1)^2C_g\overline{f}_\varepsilon,
\end{align}where $C_2(m,\chi)$ may be redefined to absorb all multiplicative factors, if necessary.

\noindent \textit{(IV) Additional collision term: }The remaining contribution from the collision operator is bounded from above in a straightforward manner. Specifically, we have
\begin{align}\label{Coll2 est}
g\overline{f}_\varepsilon \le C_g \overline{f}_\varepsilon.
\end{align}

\noindent \textit{(V) Artificial diffusion term:} To estimate the diffusion term $\Delta_{x,v}\overline{f}_\varepsilon$, we compute the spacial Laplacian of the weight function $\langle x-\chi(\varepsilon v)vt\rangle^{-m}$. A direct calculation yields that
\begin{multline*}
\Delta_x\langle x-\chi(\varepsilon v)vt\rangle^{-m} = \nabla_x\cdot[-mx\langle x-\chi(\varepsilon v)vt\rangle^{-(m+2)}]\\
= -3m\langle x-\chi(\varepsilon v)vt\rangle^{-(m+2)} +m(m+2)\langle x-\chi(\varepsilon v)vt\rangle^{-(m+4)}(x-\chi(\varepsilon v)vt),
\end{multline*}which implies the bound:
\begin{align*}
\Delta_x\langle x-\chi(\varepsilon v)vt\rangle^{-m} \le m(m+5)\langle x-\chi(\varepsilon v)vt\rangle^{-m}.
\end{align*}Combining this with  \eqref{total Delta}, we obtain
\begin{align}\label{Artifical est}
\Delta_{x,v}\overline{f}_\varepsilon \le C_2(m,\chi)(T+1)^2\overline{f}_\varepsilon,
\end{align}where the constant $C_2$ may be redefined to absorb numerical factors, if necessary. 

\noindent \textit{(VI) Artifical inhomogeneous term associated with initial data:} Since $f^\textup{in}_\varepsilon\in \mathcal{C}^2_{m,0}$, Lemma \ref{embed} yields the following estimate:
\begin{multline}\label{Initial est}
|\mathcal{L}_{\varepsilon,g}(f^\textup{in}_\varepsilon)| \le
\left(\lVert v\cdot\nabla_xf^\textup{in}_\varepsilon\rVert_{L^\infty_{m,2m}} + (C_g+1)(C_1(m) \lVert \nabla_vf^\textup{in}_\varepsilon\rVert_{L^\infty_{m,2m}}\right. \\
\left.+ C_2(m) \lVert \Delta_{v} f^\textup{in}_\varepsilon\rVert_{L^\infty_{m,2m}} + \lVert f^\textup{in}_\varepsilon\rVert_{L^\infty_{m,2m}})\right) (\langle x\rangle^{-m}\langle v\rangle^{-m})\langle v\rangle^{-m}\\
\le C(m,\chi)\lVert f^\textup{in}_\varepsilon\rVert_{\mathcal{C}^2_{m,0}}\left(C_g+1\right) (2\langle T\rangle)^m \langle x-\chi(\varepsilon v)vt\rangle^{-m}\langle v\rangle^{-m},
\end{multline}for some $C(m,\chi)>0.$

Combining \eqref{Trans est}, \eqref{Acc est}, \eqref{Coll1 est}, \eqref{Coll2 est}, \eqref{Artifical est} and \eqref{Initial est}, we assign a sufficiently large constant 
\begin{align}\label{final C'}
C' \eqdef M(m,\chi)(T+1)^2(C_g+1),
\end{align} such that the \textit{transport term} in \textit{(I)} absorbs all the rest of the terms from \textit{(II)} to \textit{(VI)} on the left-hand side of the desired inequality \eqref{4.2 desired ineq}. This allows us to conclude that
\begin{align*}
\mathcal{L}_{\varepsilon,g}(\overline{f}_\varepsilon) \ge \mathcal{L}_{\varepsilon,g}(f^\textup{in}_\varepsilon)I^\varepsilon \quad \textup{ in }\,\, \mathcal{Q}_T.
\end{align*}
By merging the upper bounds over two distinct time intervals, the desired estimate is achieved.

\textbf{Lower bound in $t\in[0,\varepsilon]$.} Since $f=f^\textup{in}_\varepsilon$ in this time interval, the desired lower bound $f\ge f^\textup{in}_\varepsilon I^\varepsilon$ trivially holds.

\textbf{Lower bound in $t\in[\varepsilon, T]$.} By comparison principle, it suffices to show 
\begin{align*}
\mathcal{L}_{\varepsilon,g}(f^\textup{in}_\varepsilon I^\varepsilon) \le \mathcal{L}_{\varepsilon,g}(f)=\mathcal{L}_{\varepsilon,g}(f^\textup{in}_\varepsilon)I^\varepsilon \quad \textup{ in }\,\, [\varepsilon, T]\times\mathbb{R}^6.
\end{align*}This inequality holds, since $$\mathcal{L}_{\varepsilon,g}(f^\textup{in}_\varepsilon I^\varepsilon) =\mathcal{L}_{\varepsilon,g}(f^\textup{in}_\varepsilon)I^\varepsilon+ f^\textup{in}_\varepsilon\partial_t I^\varepsilon, $$ and, by definition of $I^\varepsilon(t)$ in \eqref{reg interval}, we have $\partial_tI^\varepsilon \le0$ for each $t\in[\varepsilon,T]$ and  $\varepsilon>0$. Therefore, we conclude that
\begin{align*}
f^\textup{in}_\varepsilon I^\varepsilon \le f \quad \textup{in}\,\, t\ge\varepsilon.
\end{align*}

By combining the lower bounds on the two distinct time intervals, we obtain the desired estimate.

In the above proof, the constant $M$ was written as if it depends on $\chi$. More precisely, increasing $\|\chi\|_{C^{2}}$ increases the corresponding value $M(\chi)$. However, this dependence on $\chi$ can in fact be removed. We may define
$$
    M(m) \coloneqq \inf\{\, M(m,\chi) : \chi\in C_c^2(\mathbb{R}^3) \textup{ and satisfies }\eqref{def cutoff}\,\},
$$
which yields a choice of $M$ independent of $\chi$. The infimum $M(m)>0$ exists for each fixed $m$. This completes the proof.
\end{proof}

\end{document}
